% GENERATED FILE. Edit canonical lessons, then run npm run generate:books.
% canonical-source-hash: fnv1a64:e1e5c130be54657c
% canonical-lessons: SA-C158-lekhani, SA-C158-masi, SA-C158-pattika, SA-C158-dipika, SA-C158-kapatam, SA-R158-first-pass-words-from-chapters-150-153, SA-R158-second-pass-words-from-chapters-154-158

\chapter{Pen, Ink, Slate, Small lamp, Door panel}
\label{ch:sa-pen-ink-slate-small-lamp-door-panel}
\hlchaptermodality{158}

\begin{chapteropening}
\textbf{By the end of this chapter:} \emph{I can say a pen, ink, a slate, a small lamp and a door panel.}

Complete the last lesson of chapter 158: I can say a pen, ink, a slate, a small lamp and a door panel.
\end{chapteropening}

\section[\texorpdfstring{\sk{लेखनी} (lekhanī)}{लेखनी (lekhanī)}]{\sk{लेखनी} (lekhanī) — a pen}

\label{lesson:SA-C158-lekhani}



\begin{quote}
Before the new one: say the Sanskrit for tired, then the Sanskrit for hungry.
\end{quote}

\subsection*{You'll want to know: \sk{लेखनी}}
\textbf{\sk{लेखनी}} — \emph{lekhanī} — \textquotedblleft{}a pen\textquotedblright{}.

\begin{grammarlens}[title={What you've built}]
One more word for this chapter.
\end{grammarlens}

\subsection*{Your turn}
\begin{itemize}
\raggedright
  \item \emph{Say it:} \emph{lekhanī}
  \item \emph{Say it:} \emph{lekhanī}, once more
  \item \emph{Say it:} say \emph{lekhanī}
  \item \emph{Recall:} say the Sanskrit for tired, then the Sanskrit for hungry, then say \emph{lekhanī} again
\end{itemize}

\begin{tcolorbox}[breakable,colback=teal!4,colframe=teal!35!black,title={Before you move on}]
What does \sk{लेखनी} mean? (\textquotedblleft{}A pen\textquotedblright{}.) Say it once more.
\end{tcolorbox}
\section[\texorpdfstring{\sk{मसी} (masī)}{मसी (masī)}]{\sk{मसी} (masī) — ink}

\label{lesson:SA-C158-masi}



\begin{quote}
Before the new one: say the Sanskrit for hungry, then the Sanskrit for a pen.
\end{quote}

\subsection*{You'll want to know: \sk{मसी}}
\textbf{\sk{मसी}} — \emph{masī} — \textquotedblleft{}ink\textquotedblright{}.

\begin{grammarlens}[title={What you've built}]
One more word for this chapter.
\end{grammarlens}

\subsection*{Your turn}
\begin{itemize}
\raggedright
  \item \emph{Say it:} \emph{masī}
  \item \emph{Say it:} \emph{masī}, once more
  \item \emph{Say it:} say \emph{masī}
  \item \emph{Recall:} say the Sanskrit for hungry, then the Sanskrit for a pen, then say \emph{masī} again
\end{itemize}

\begin{tcolorbox}[breakable,colback=teal!4,colframe=teal!35!black,title={Before you move on}]
What does \sk{मसी} mean? (\textquotedblleft{}Ink\textquotedblright{}.) Say it once more.
\end{tcolorbox}
\section[\texorpdfstring{\sk{पट्टिका} (paṭṭikā)}{पट्टिका (paṭṭikā)}]{\sk{पट्टिका} (paṭṭikā) — a slate, a tablet}

\label{lesson:SA-C158-pattika}



\begin{quote}
Before the new one: say the Sanskrit for a pen, then the Sanskrit for ink.
\end{quote}

\subsection*{You'll want to know: \sk{पट्टिका}}
\textbf{\sk{पट्टिका}} — \emph{paṭṭikā} — \textquotedblleft{}a slate, a tablet\textquotedblright{}.

\begin{grammarlens}[title={What you've built}]
One more word for this chapter.
\end{grammarlens}

\subsection*{Your turn}
\begin{itemize}
\raggedright
  \item \emph{Say it:} \emph{paṭṭikā}
  \item \emph{Say it:} \emph{paṭṭikā}, once more
  \item \emph{Say it:} say \emph{paṭṭikā}
  \item \emph{Recall:} say the Sanskrit for a pen, then the Sanskrit for ink, then say \emph{paṭṭikā} again
\end{itemize}

\begin{tcolorbox}[breakable,colback=teal!4,colframe=teal!35!black,title={Before you move on}]
What does \sk{पट्टिका} mean? (\textquotedblleft{}A slate, a tablet\textquotedblright{}.) Say it once more.
\end{tcolorbox}
\section[\texorpdfstring{\sk{दीपिका} (dīpikā)}{दीपिका (dīpikā)}]{\sk{दीपिका} (dīpikā) — a small lamp}

\label{lesson:SA-C158-dipika}



\begin{quote}
Before the new one: say the Sanskrit for ink, then the Sanskrit for a slate.
\end{quote}

\subsection*{You'll want to know: \sk{दीपिका}}
\textbf{\sk{दीपिका}} — \emph{dīpikā} — \textquotedblleft{}a small lamp\textquotedblright{}.

\begin{grammarlens}[title={What you've built}]
One more word for this chapter.
\end{grammarlens}

\subsection*{Your turn}
\begin{itemize}
\raggedright
  \item \emph{Say it:} \emph{dīpikā}
  \item \emph{Say it:} \emph{dīpikā}, once more
  \item \emph{Say it:} say \emph{dīpikā}
  \item \emph{Recall:} say the Sanskrit for ink, then the Sanskrit for a slate, then say \emph{dīpikā} again
\end{itemize}

\begin{tcolorbox}[breakable,colback=teal!4,colframe=teal!35!black,title={Before you move on}]
What does \sk{दीपिका} mean? (\textquotedblleft{}A small lamp\textquotedblright{}.) Say it once more.
\end{tcolorbox}
\section[\texorpdfstring{\sk{कपाटम्} (kapāṭam)}{कपाटम् (kapāṭam)}]{\sk{कपाटम्} (kapāṭam) — a door panel, a shutter}

\label{lesson:SA-C158-kapatam}



\begin{quote}
Before the new one: say the Sanskrit for a slate, then the Sanskrit for a small lamp.
\end{quote}

\subsection*{You'll want to know: \sk{कपाटम्}}
\textbf{\sk{कपाटम्}} — \emph{kapāṭam} — \textquotedblleft{}a door panel, a shutter\textquotedblright{}.

\begin{grammarlens}[title={What you've built}]
One more word for this chapter.
\end{grammarlens}

\subsection*{Your turn}
\begin{itemize}
\raggedright
  \item \emph{Say it:} \emph{kapāṭam}
  \item \emph{Say it:} \emph{kapāṭam}, once more
  \item \emph{Say it:} say \emph{kapāṭam}
  \item \emph{Recall:} say the Sanskrit for a slate, then the Sanskrit for a small lamp, then say \emph{kapāṭam} again
\end{itemize}

\begin{tcolorbox}[breakable,colback=teal!4,colframe=teal!35!black,title={Before you move on}]
What does \sk{कपाटम्} mean? (\textquotedblleft{}A door panel, a shutter\textquotedblright{}.) Say it once more.
\end{tcolorbox}
\section[(dialogue)]{Words from chapters 150-153 — a first pass}

\label{lesson:SA-R158-first-pass-words-from-chapters-150-153}



\begin{quote}
Say the words for a small lamp and a door panel.
\end{quote}

\subsection*{Your turn}
Say these aloud:

\begin{itemize}
\raggedright
  \item a remedy, a wound, a swelling, vomiting, a stomach ache
  \item enthusiasm, courage, laziness, jealousy, pride
  \item a period, a turn, a lunar day, a calendar date, a century
  \item from here, from there, on both sides, all around, on all sides
\end{itemize}

\begin{tcolorbox}[breakable,colback=teal!4,colframe=teal!35!black,title={Before you move on}]
A remedy? (\textbf{\sk{भेषजम्}}, \emph{bheṣajam}.) A wound? (\textbf{\sk{व्रणः}}, \emph{vraṇaḥ}.) A swelling? (\textbf{\sk{शोथः}}, \emph{śothaḥ}.) Vomiting? (\textbf{\sk{वमनम्}}, \emph{vamanam}.) A stomach ache? (\textbf{\sk{उदरपीडा}}, \emph{udarapīḍā}.) Enthusiasm? (\textbf{\sk{उत्साहः}}, \emph{utsāhaḥ}.) Courage? (\textbf{\sk{धैर्यम्}}, \emph{dhairyam}.) Laziness? (\textbf{\sk{आलस्यम्}}, \emph{ālasyam}.) Jealousy? (\textbf{\sk{मत्सरः}}, \emph{matsaraḥ}.) Pride? (\textbf{\sk{गर्वः}}, \emph{garvaḥ}.) A period? (\textbf{\sk{कालः}}, \emph{kālaḥ}.) A turn? (\textbf{\sk{वारः}}, \emph{vāraḥ}.) A lunar day? (\textbf{\sk{तिथिः}}, \emph{tithiḥ}.) A calendar date? (\textbf{\sk{दिनांकः}}, \emph{dināṅkaḥ}.) A century? (\textbf{\sk{शताब्दी}}, \emph{śatābdī}.) From here? (\textbf{\sk{इतः}}, \emph{itaḥ}.) From there? (\textbf{\sk{ततः}}, \emph{tataḥ}.) On both sides? (\textbf{\sk{उभयतः}}, \emph{ubhayataḥ}.) All around? (\textbf{\sk{परितः}}, \emph{paritaḥ}.) On all sides? (\textbf{\sk{समन्ततः}}, \emph{samantataḥ}.)
\end{tcolorbox}
\section[(dialogue)]{Words from chapters 154-158 — a second pass}

\label{lesson:SA-R158-second-pass-words-from-chapters-154-158}



\begin{quote}
Say the words for if and provided that.
\end{quote}

\subsection*{Your turn}
Say these aloud:

\begin{itemize}
\raggedright
  \item if, provided that, even so, although, or else
  \item fresh, worn out, hot, chilly, harsh
  \item dear, unpleasant, calm, fierce, cruel
  \item angry, afraid, startled, tired, hungry
  \item a pen, ink, a slate, a small lamp, a door panel
\end{itemize}

\begin{tcolorbox}[breakable,colback=teal!4,colframe=teal!35!black,title={Before you move on}]
If? (\textbf{\sk{यदि}}, \emph{yadi}.) Provided that? (\textbf{\sk{चेत्}}, \emph{cet}.) Even so? (\textbf{\sk{तथापि}}, \emph{tathāpi}.) Although? (\textbf{\sk{यद्यपि}}, \emph{yadyapi}.) Or else? (\textbf{\sk{अथवा}}, \emph{athavā}.) Fresh? (\textbf{\sk{नूतनः}}, \emph{nūtanaḥ}.) Worn out? (\textbf{\sk{जीर्णः}}, \emph{jīrṇaḥ}.) Hot? (\textbf{\sk{उष्णः}}, \emph{uṣṇaḥ}.) Chilly? (\textbf{\sk{शीतः}}, \emph{śītaḥ}.) Harsh? (\textbf{\sk{कठोरः}}, \emph{kaṭhoraḥ}.) Dear? (\textbf{\sk{प्रियः}}, \emph{priyaḥ}.) Unpleasant? (\textbf{\sk{अप्रियः}}, \emph{apriyaḥ}.) Calm? (\textbf{\sk{शान्तः}}, \emph{śāntaḥ}.) Fierce? (\textbf{\sk{चण्डः}}, \emph{caṇḍaḥ}.) Cruel? (\textbf{\sk{क्रूरः}}, \emph{krūraḥ}.) Angry? (\textbf{\sk{कुपितः}}, \emph{kupitaḥ}.) Afraid? (\textbf{\sk{भीतः}}, \emph{bhītaḥ}.) Startled? (\textbf{\sk{चकितः}}, \emph{cakitaḥ}.) Tired? (\textbf{\sk{श्रान्तः}}, \emph{śrāntaḥ}.) Hungry? (\textbf{\sk{क्षुधितः}}, \emph{kṣudhitaḥ}.) A pen? (\textbf{\sk{लेखनी}}, \emph{lekhanī}.) Ink? (\textbf{\sk{मसी}}, \emph{masī}.) A slate? (\textbf{\sk{पट्टिका}}, \emph{paṭṭikā}.) A small lamp? (\textbf{\sk{दीपिका}}, \emph{dīpikā}.) A door panel? (\textbf{\sk{कपाटम्}}, \emph{kapāṭam}.)
\end{tcolorbox}
